\documentclass[11pt,letterpaper]{article}
\usepackage[T1]{fontenc}
\usepackage[utf8]{inputenc}
\usepackage{lmodern}
\usepackage[margin=1in,headheight=15pt]{geometry}
\usepackage{amsmath,amssymb,amsthm,mathtools,mathrsfs,aliascnt}
\usepackage{microtype,booktabs,tabularx,array,enumitem}
\usepackage{xcolor,fancyhdr,titlesec,needspace,etoolbox}
\definecolor{Ink}{HTML}{243747}
\definecolor{Accent}{HTML}{37666D}
\usepackage[colorlinks=true,linkcolor=Ink,citecolor=Accent,urlcolor=Accent]{hyperref}
\usepackage{xurl}
\usepackage[nameinlink,noabbrev]{cleveref}
\titleformat{\section}{\Large\bfseries\color{Ink}}{\thesection}{.65em}{}
\titleformat{\subsection}{\large\bfseries\color{Accent}}{\thesubsection}{.65em}{}
\pretocmd{\section}{\Needspace{8\baselineskip}}{}{}
\pretocmd{\subsection}{\Needspace{4\baselineskip}}{}{}
\pagestyle{fancy}\fancyhf{}
\fancyhead[L]{\small Three duals at surreal scales}
\fancyhead[R]{\small Strong sums and completion defects}
\fancyfoot[C]{\thepage}
\setlength{\parskip}{.28em}
\setlength{\parindent}{1.15em}
\setlength{\emergencystretch}{3em}
\widowpenalty=10000\clubpenalty=10000
\setlist{itemsep=.2em,topsep=.35em,leftmargin=2em}
\setcounter{tocdepth}{2}
\numberwithin{equation}{section}
\newtheorem{theorem}{Theorem}[section]
\newcommand{\newaliastheorem}[3]{%
  \newaliascnt{#1}{theorem}\newtheorem{#1}[#1]{#2}%
  \aliascntresetthe{#1}\crefname{#1}{#2}{#3}}
\crefname{theorem}{Theorem}{Theorems}
\newaliastheorem{proposition}{Proposition}{Propositions}
\newaliastheorem{lemma}{Lemma}{Lemmas}
\newaliastheorem{corollary}{Corollary}{Corollaries}
\theoremstyle{definition}
\newaliastheorem{definition}{Definition}{Definitions}
\newaliastheorem{example}{Example}{Examples}
\theoremstyle{remark}
\newaliastheorem{remark}{Remark}{Remarks}
\newaliastheorem{warning}{Warning}{Warnings}
\newcommand{\C}{\mathbb C}
\newcommand{\R}{\mathbb R}
\newcommand{\Q}{\mathbb Q}
\newcommand{\Z}{\mathbb Z}
\newcommand{\N}{\mathbb N}
\newcommand{\No}{\mathbf{No}}
\newcommand{\K}{K}
\newcommand{\F}{F}
\newcommand{\HG}{\mathcal H_\Gamma}
\newcommand{\VG}{V_\Gamma}
\newcommand{\WG}{W_\Gamma}
\newcommand{\EG}{E_\Gamma(V)}
\newcommand{\CG}{C_\Gamma(V)}
\newcommand{\supp}{\operatorname{supp}}
\newcommand{\Span}{\operatorname{span}}
\newcommand{\Hom}{\operatorname{Hom}}
\newcommand{\End}{\operatorname{End}}
\newcommand{\res}{\operatorname{res}}
\newcommand{\rank}{\operatorname{rank}}
\newcommand{\st}{\operatorname{st}}
\newcommand{\Dstr}{D_{\mathrm{str}}}
\newcommand{\Dcont}{D_{\mathrm{cont}}}
\newcommand{\Drep}{D_{\mathrm{rep}}}
\newcommand{\sH}{\mathop{\sum}\nolimits^{\mathrm H}}
\newcommand{\abs}[1]{\lvert#1\rvert}
\newcommand{\norm}[1]{\lVert#1\rVert}
\newcommand{\ip}[2]{\langle#1,#2\rangle}
\newcommand{\cyc}{\Gamma\simeq_{\mathrm{ord}}\Z}
\newcommand{\Commit}{048b72cf7cbfc8ab246e4f73788c10460cb3f6e0}
\newcolumntype{Y}{>{\raggedright\arraybackslash}X}
\hypersetup{pdftitle={Three Duals at Surreal Scales: Strong Hahn Operators, Completion Defects, and Invisible Functionals},pdfauthor={Research manuscript prepared for Vladimir Reshetnikov},pdfsubject={Hahn vector spaces, continuous and strong duality, completion and scalar enlargement}}
\begin{document}
\hypersetup{pageanchor=false}
\begin{titlepage}
\centering
\vspace*{.10in}
{\large\scshape A research article for the Surreal project\par}
\vspace{.20in}
{\Huge\bfseries\color{Ink} Three Duals at\\[.15em] Surreal Scales\par}
\vspace{.15in}
{\LARGE Strong Hahn Operators, Completion Defects,\\[.15em] and Invisible Functionals\par}
\vspace{.20in}
{\large September 22, 2026\par}
\vspace{.20in}
\begin{minipage}{.94\textwidth}
\begin{abstract}\noindent
Let $K=k((t^\Gamma))$ and $V_\Gamma=V((t^\Gamma))$, where $V$ is a
$k$-vector space and $\Gamma$ is a set-sized ordered abelian group. We
separate the algebraic scalar extension, its valuation completion, and its
closure under strong Hahn summation. The completion consists exactly of
those series whose coefficients below each valuation cut span a
finite-dimensional $k$-space. In infinite coefficient dimension it equals
the full Hahn space exactly when the nonzero group $\Gamma$ is ordered
isomorphic to $\mathbb Z$. Under a cofinal exponent enlargement this
completion criterion is preserved; under a noncofinal enlargement its
intersection with the original Hahn space collapses to the algebraic
scalar extension.

We also classify every strong $K$-linear operator as a single Hahn series
of arbitrary $k$-linear coefficient maps. The continuous dual has an exact
restriction sequence: its image consists of pointwise Hahn-valued
coefficient functionals with a common lower support bound, while its
kernel consists of continuous functionals vanishing on every constant
vector. These latter functionals separate precisely the vectors outside
the completed scalar extension. For every infinite-dimensional $V$ and
noncyclic $\Gamma\ne0$, there are two distinct failures of strongness:
continuous maps can destroy summability, or preserve the summability of
a particular monomial family but fail to commute with its sum.

For infinite-dimensional complex Hilbert $H$ and noncyclic $\Gamma$,
this produces three strictly nested duals: inner-product-represented, strong, and valuation-continuous.
The first gap is exactly the Hahn extension of the ordinary algebraic-dual
quotient $H^\#/H'$. All strong maps are bounded by a positive real Hahn
scalar, even when their coefficient maps are discontinuous in the
ordinary Hilbert norm. An explicit choice-based hyperplane has a strong
bounded projection but zero orthogonal complement, and the distance from
an exterior constant vector has no infimum in the scalar field.
\end{abstract}
\end{minipage}
\vfill
\begin{minipage}{.94\textwidth}\small
\textbf{Scope and status.} Complete mathematical proofs are supplied for
the stated results. The proposed contribution is the combined
completion--duality--scalar-enlargement classification, not the invention
of Hahn summation, non-Archimedean Hahn--Banach, or failures of Riesz
representation. No named published conjecture is claimed solved. A
targeted repository and primary-literature comparison did not locate this
exact package; this is not a certification of priority. The proofs have
not been independently refereed or formally verified. All workspaces and
choices in the proofs are set-sized.
\end{minipage}
\end{titlepage}
\hypersetup{pageanchor=true}
\pagenumbering{roman}\tableofcontents\clearpage
\pagenumbering{arabic}

\section{The questions and the contribution}\label{sec:intro}

\subsection{Why continuity is not the same as respecting a formal sum}
An infinite-dimensional extension of complex analysis to surcomplex
scalars has to specify more than its scalar field. It must specify its
vectors, the admissible infinite sums, and the meaning of continuity.
For the natural coefficient space
\[
 H((t^\Gamma))=\left\{\sH_{\gamma\in\Gamma}h_\gamma t^\gamma:
 h_\gamma\in H,\ \supp h\text{ is well ordered}\right\},
\]
three plausible requirements on a scalar-valued linear functional are
not interchangeable. It may be given by an inner product. It may commute
with every admissible Hahn sum. Or it may merely be continuous when
closeness means agreement below a sufficiently high valuation cut.

The issue is visible before introducing a Hilbert norm. Given an
ordinary vector space $V$, form the algebraic scalar extension
$K\otimes_k V$. It embeds in $V((t^\Gamma))$, and its vectors have
coefficients in one finite-dimensional subspace of $V$. Completing this
subspace need not give all of $V((t^\Gamma))$. A Hahn sum can contain
infinitely many linearly independent coefficients below a single
valuation cut. Such a vector is permitted by well-ordered support, but
no finite coefficient span can match it to that precision.

This distinction leads to four concrete questions. Which Hahn vectors
belong to the completed algebraic scalar extension? Which linear maps
respect all strong sums? What can a continuous functional detect that
its values on constants do not determine? What happens to these answers
when additional surreal scales are admitted? The theorems below give
exact answers, including the exceptional cases.

\subsection{Main statements at a glance}
For a nonzero ordered abelian group $\Gamma$, write
$V_\Gamma=V((t^\Gamma))$, $K=k((t^\Gamma))$,
\[
 E_\Gamma(V)=\Span_K V,\qquad
 C_\Gamma(V)=\overline{E_\Gamma(V)}^{\,v}.
\]
Our convention is that a strong map both preserves admissibility of all
strongly summable families and commutes with their sums.

\begin{theorem}[Completion and enlargement; proved in
\cref{sec:completion,sec:enlargement}]\label{thm:overview-completion}
For every $x\in V_\Gamma$,
\begin{align}
 x\in E_\Gamma(V)
 &\iff \dim_k\Span\{x_\delta:\delta\in\Gamma\}<\infty,\label{eq:overview-E}\\
 x\in C_\Gamma(V)
 &\iff (\forall\gamma\in\Gamma)\ 
 \dim_k\Span\{x_\delta:\delta<\gamma\}<\infty.\label{eq:overview-C}
\end{align}
If $V$ is infinite-dimensional, then
$C_\Gamma(V)=V_\Gamma$ exactly when $\cyc$.
For an ordered subgroup inclusion $\Gamma\subseteq\Delta$,
\[
 C_\Delta(V)\cap V_\Gamma=
 \begin{cases}
 C_\Gamma(V),&\Gamma\text{ is cofinal in }\Delta,\\
 E_\Gamma(V),&\Gamma\text{ is not cofinal in }\Delta.
 \end{cases}
\]
\end{theorem}

\begin{theorem}[Strong maps and the continuous dual; proved in
\cref{sec:strong,sec:continuous}]\label{thm:overview-duals}
For arbitrary $k$-vector spaces $V,W$,
\[
 \Hom_K^{\mathrm{str}}(V_\Gamma,W_\Gamma)
 \simeq \Hom_k(V,W)((t^\Gamma)).
\]
Let $P_\Gamma(V)$ be the space of $k$-linear maps $A:V\to K$ for which
$v(Au)\ge s$ for all $u\in V$, with some common $s\in\Gamma$.
There is an exact sequence
\[
 0\longrightarrow N_\Gamma(V)\longrightarrow
 \Hom_K^{\mathrm{cont}}(V_\Gamma,K)
 \xrightarrow{\ \res\ } P_\Gamma(V)\longrightarrow0,
\]
where $N_\Gamma(V)$ consists of the continuous functionals zero on $V$.
Their common kernel is exactly $C_\Gamma(V)$.
\end{theorem}

For infinite-dimensional $V$ and noncyclic $\Gamma$, the kernel
$N_\Gamma(V)$ is nonzero. Moreover, some restrictions in $P_\Gamma(V)$
are not the restrictions of any strong map: their individual values are
Hahn series, but their \emph{global} coefficient support is not well
ordered. Thus passing from strong maps to continuous maps adds two
separate kinds of freedom, not just an undetectable additive error.

For a complex Hilbert coefficient space the additional comparison with
inner products is
\begin{equation}\label{eq:overview-three}
 \Drep\ \subseteq\ \Dstr\ \subseteq\ \Dcont,
 \qquad
 \Dstr/\Drep\simeq (H^\#/H')((t^\Gamma)).
\end{equation}
Both inclusions are strict when $H$ is infinite-dimensional and
$\Gamma$ is noncyclic. When $\Gamma\simeq_{\mathrm{ord}}\Z$, only the
first is strict. In finite coefficient dimension all three coincide.

\subsection{What has already been done, and what is not claimed}
The repository's \emph{Infinite-Dimensional Hahn Spectral Theory}
already constructs the coefficient-convolution inner product and proves
that every everywhere-defined adjointable operator is a Hahn series of
ordinary bounded operators. It also proves spectral and closed-range
phenomena for that space \cite{RepoSpectral}. We do not present those
results as new. Our operator theorem replaces the adjoint hypothesis by
preservation of strong sums, permits \emph{arbitrary algebraic}
coefficient maps, and is then used to compare with the strictly larger
continuous dual. The coefficient-rank completion criterion and its
cofinal/noncofinal enlargement dichotomy address a different question
from the repository's persistence of spectral range defects.

Formal summability spaces, strongly linear maps, and the resulting
operator algebras are established notions. In particular,
Bagayoko--Krapp--Kuhlmann--Panazzolo--Serra give a general axiomatic
framework and examples of nonstrong linear maps
\cite{BKKPS}. Our definition follows that usage. Their scalar-field
linearity conventions must not be conflated with ours: the main operator
classification here is \emph{linear over the entire Hahn field} $K$,
not merely over its coefficient field $k$.

The extension argument belongs to the classical non-Archimedean
Hahn--Banach tradition initiated by Ingleton \cite{Ingleton}; Morillon
states its one-step form and discusses its choice requirements
\cite{Morillon}. We give the needed arbitrary-ordered-group proof in
full. Spherical completeness and extension are tools, not claimed
innovations. Non-Archimedean failures of orthogonal complementation
are likewise established; Aguayo--Nova study them in a $c_0$-based
model \cite{AN}. Our final hyperplane theorem specifies a different
coefficient-Hilbert model and an exact missing-infimum cut.

The proposed originality lies in the precise classification package,
especially \cref{thm:completion,thm:cyclic,thm:basechange,thm:restriction,thm:separation,thm:twofailures}. The inspected repository inventory and
the primary sources consulted are recorded in \cref{sec:audit} and the
accompanying audit file. This comparison is deliberately narrower than
an assertion that no equivalent theorem exists anywhere in print.

\section{Hahn vector spaces and the two notions of infinity}\label{sec:setup}

\subsection{Supports, scalar multiplication, and strong families}
Fix a field $k$ and a set-sized ordered abelian group $\Gamma$, written
additively. Its order is translation invariant. For a $k$-vector space
$V$, a Hahn vector is a function $x:\Gamma\to V$ whose nonzero support
is well ordered in the given order. We use formal notation
\[
 x=\sH_{\gamma\in\Gamma}x_\gamma t^\gamma,
 \qquad V_\Gamma=V((t^\Gamma)).
\]
The zero vector has empty support. The scalar field is
$K=k((t^\Gamma))$. Its multiplication and its action on $V_\Gamma$
are coefficient convolutions. In particular,
\[
 (a x)_\eta=\sum_{\alpha+\beta=\eta}a_\alpha x_\beta.
\]
The superscript H emphasizes formal Hahn summation; it never means
ordinary summation of infinitely many coefficients at one exponent.

\begin{definition}\label{def:strong-family}
A set-indexed family $(x_j)_{j\in J}$ in $V_\Gamma$ is strongly
Hahn summable if $\bigcup_j\supp x_j$ is well ordered and, for every
$\gamma$, only finitely many $j$ have $(x_j)_\gamma\ne0$.
Its sum is defined by the finite coefficient sums
$[t^\gamma](\sH_j x_j)=\sum_j(x_j)_\gamma$.
\end{definition}

Every Hahn vector is the strong sum of its monomial coefficient vectors.
An infinite family of nonzero constants is not strong. Even an ordinary
convergent numerical series, such as $\sum_{n\ge1}2^{-n}$, is not
that family's Hahn sum: the exponent zero has infinitely many
contributors. The familiar numerical value may be used as a
\emph{single coefficient}, which is a different operation.

\begin{lemma}[Classical support facts]\label{lem:support}
If $A,B\subseteq\Gamma$ are well ordered, then $A+B$ is well ordered,
and each $\eta$ has only finitely many representations
$\eta=\alpha+\beta$ with $\alpha\in A$, $\beta\in B$.
A finite union of well-ordered subsets is well ordered.
If $S\subseteq\Gamma_{>0}$ is well ordered, the monoid of finite sums
of elements of $S$ is well ordered; for each exponent there are only
finitely many words in $S$ having that sum, allowing all finite lengths.
\end{lemma}
\begin{proof}
For the two-set statement, every sequence in a well-ordered set has an
infinite nondecreasing subsequence: successively choose the minimum of
each remaining tail. A decreasing sequence of sums, after passing to
such a subsequence of its $A$-coordinates, would force a decreasing
sequence in $B$. Infinitely many representations of one sum similarly
give strictly increasing $A$-coordinates and strictly decreasing
$B$-coordinates. The finite-union assertion follows by selecting a
minimum in each nonempty intersection. The positive-monoid assertion is
the classical Hahn--Neumann support lemma \cite{Neumann}; it is needed
below only for the standard binomial square-root construction. The
operator and completion classifications use the two-set assertion, whose
proof was just given.
\end{proof}

Strong sums may be regrouped along any map of index sets. Indeed,
regrouping can only reduce support, and each coefficient involves
finitely many original indices. Scalar multiplication commutes with
strong sums by \cref{lem:support}. These statements supply all
infinite rearrangements used in the operator proofs.

\subsection{Valuation topology and bounded shifts}
Define
\[
 v(x)=\min\supp x\quad(x\ne0),\qquad v(0)=+\infty.
\]
For $a\in K^\times$,
$v(ax)=v(a)+v(x)$, and
$v(x+y)\ge\min\{v(x),v(y)\}$. The equality for scalar multiplication
holds because the product's first coefficient is a nonzero scalar
multiple of the first vector coefficient. It entails no norm on $V$.

Except where expressly stated, assume $\Gamma\ne0$. The intrinsic
valuation topology has a neighborhood basis
\[
 U_\gamma(x)=\{y:v(y-x)>\gamma\},\qquad\gamma\in\Gamma.
\]
Closed valuation balls
$B(x,\gamma)=\{y:v(y-x)\ge\gamma\}$ give equivalent precision tests:
one can replace $\gamma$ by $\gamma+\eta$, with any fixed $\eta>0$.
The topology is Hausdorff, since a nonzero difference has a finite
valuation and can be excluded by a higher precision.

\begin{lemma}[Continuity is a uniform valuation-shift bound]
\label{lem:continuous-shift}
A $K$-linear map $L:V_\Gamma\to W_\Gamma$ is valuation-continuous
if and only if there exists $s\in\Gamma$ such that
\begin{equation}\label{eq:shift}
 v(Lx)\ge v(x)+s\qquad(x\in V_\Gamma).
\end{equation}
The same assertion holds for a $K$-linear map on any subspace with the
induced valuation topology.
\end{lemma}
\begin{proof}
A bound of the displayed form proves continuity at zero, hence
everywhere. Conversely, continuity at zero supplies $\delta\in\Gamma$
such that $v(z)>\delta$ implies $v(Lz)>0$. Fix $\eta>0$. For nonzero
$x$, put $a=t^{\delta+\eta-v(x)}$. Then $v(ax)=\delta+\eta$, so
\[
 v(Lx)>v(x)-\delta-\eta.
\]
Thus $s=-\delta-\eta$ gives the required weak inequality. The argument
uses only scalar stability of the domain, and therefore works on a
subspace as well.
\end{proof}

\begin{warning}[Constants carry the discrete valuation topology]
If $V$ is originally normed, its embedded constant copy in $V_\Gamma$
does not carry that norm topology: every nonzero constant has valuation
zero. A map discontinuous in the ordinary norm can therefore become
continuous after coefficientwise Hahn extension. This is not a
contradiction to ordinary functional analysis.
\end{warning}

\section{All strong linear maps have one global operator support}
\label{sec:strong}

\subsection{The exact Hom theorem}
\begin{definition}
A $K$-linear map $T:V_\Gamma\to W_\Gamma$ is strong if, for every
strongly Hahn-summable family $(x_j)$, the family $(Tx_j)$ is strongly
Hahn summable and $T(\sH_jx_j)=\sH_jTx_j$.
\end{definition}
Both clauses are part of the definition. In \cref{sec:twofailures}
continuous functionals will exhibit two different obstructions.

\begin{theorem}[Global-support classification]\label{thm:strong}
For all $k$-vector spaces $V,W$ and every ordered abelian group
$\Gamma$, including the zero group, the convolution map gives a
canonical $K$-linear isomorphism
\begin{equation}\label{eq:strong-Hom}
 \Hom_k(V,W)((t^\Gamma))
 \xrightarrow{\ \sim\ }
 \Hom_K^{\mathrm{str}}(V_\Gamma,W_\Gamma).
\end{equation}
Explicitly, a Hahn series $A=\sH_\gamma A_\gamma t^\gamma$ acts by
\begin{equation}\label{eq:operator-convolution}
 (A x)_\eta=\sum_{\gamma+\delta=\eta}A_\gamma(x_\delta).
\end{equation}
The set $\{\gamma:A_\gamma\ne0\}$ is one well-ordered support,
independent of $x$. No topology on $V$ or $W$ is assumed.
\end{theorem}
\begin{proof}
\emph{Necessity of one global support.}
Let $T$ be strong. For a constant $u\in V$ define
$A_\gamma(u)=[t^\gamma]Tu$. Coefficient extraction shows that every
$A_\gamma$ is $k$-linear. Put $S=\{\gamma:A_\gamma\ne0\}$.
If $S$ were not well ordered, choice would provide a strictly decreasing
sequence $\gamma_1>\gamma_2>\cdots$ in $S$. Select $u_n\in V$ with
$A_{\gamma_n}(u_n)\ne0$.

The family $x_n=t^{-\gamma_n}u_n$ is strong: its singleton exponents
form a strictly increasing sequence, and every exponent occurs only
once. But $K$-linearity gives $Tx_n=t^{-\gamma_n}Tu_n$, whose coefficient
at exponent zero is $A_{\gamma_n}(u_n)\ne0$ for every $n$. Its images
therefore violate the finite-fiber condition. This contradicts
strongness and proves that $S$ is well ordered.

\emph{Determination by constants.}
Every $x$ is the strong sum of $t^\delta x_\delta$ over its support.
Strongness and $K$-linearity yield
\[
 Tx=\sH_\delta t^\delta T(x_\delta)
   =\sH_{\delta,\gamma}A_\gamma(x_\delta)t^{\delta+\gamma}.
\]
The common support $S$ and \cref{lem:support} justify the final
expansion and make its coefficient formula exactly
\eqref{eq:operator-convolution}. Testing constants also proves
uniqueness of the coefficient maps.

\emph{Sufficiency.}
Start instead with a well-ordered support of maps $A_\gamma$.
For one input $x$, \cref{lem:support} makes
\eqref{eq:operator-convolution} a Hahn vector. It commutes with scalar
multiplication, by finite rearrangement at each coefficient, and is
additive. If $(x_j)$ is strong, let $B=\bigcup_j\supp x_j$.
All image supports lie in $S+B$. At a fixed exponent there are finitely
many pairs $(\gamma,\delta)$, and for each such $\delta$ finitely many
$j$ contribute. Thus the image family is strong, and the same finite
rearrangement proves commutation with its sum.
\end{proof}

The descending-support test is an automatic uniformity principle: an
assumption about all summable families rules out a global support defect
that may not be visible on any one constant vector. It does not use
Baire category, completeness of $V$, or an adjoint.

\begin{corollary}[Operator calculus and automatic continuity]
\label{cor:strong-calculus}
Composition corresponds to coefficient convolution. In particular,
\[
 \End_K^{\mathrm{str}}(V_\Gamma)
 \simeq \End_k(V)((t^\Gamma))
\]
as unital algebras. If $A\ne0$ and $s=\min\supp A$, then
$v(Ax)\ge v(x)+s$ for all $x$. Every strong map is valuation-continuous.
\end{corollary}
\begin{proof}
For composable $A,B$, the coefficient of $B\circ A$ at $\eta$ is
$\sum_{\beta+\alpha=\eta}B_\beta A_\alpha$. All sums are finite by
\cref{lem:support}, and the action on a vector agrees after regrouping.
For the estimate, every nonzero output coefficient lies at a sum of an
exponent at least $s$ and one at least $v(x)$. Apply
\cref{lem:continuous-shift}.
\end{proof}

\subsection{Pointwise supports alone do not suffice}
For a mere $k$-linear map $A:V\to K$, every individual $A(u)$ has a
well-ordered support. It need not follow that
\[
 \bigcup_{u\in V}\supp A(u)
\]
is well ordered. For instance, take a Hamel basis containing distinct
$e_n$, a strictly decreasing positive sequence $\delta_n$, and set
$A(e_n)=t^{\delta_n}$, with zero values on all remaining basis vectors.
Every $A(u)$ has finite support. The global support contains the
entire decreasing sequence. This example cannot be inserted in
\eqref{eq:operator-convolution} on arbitrary Hahn vectors. It can,
however, be extended to a continuous functional by a different,
noncanonical construction. That is one of the distinctions explained
in \cref{sec:continuous}.

\section{The completion is a coefficient-rank condition}
\label{sec:completion}

\subsection{The algebraic scalar extension inside the Hahn space}
The constant embedding $V\hookrightarrow V_\Gamma$ induces a map
$K\otimes_kV\to V_\Gamma$. If $u_1,\ldots,u_d$ are linearly independent
in $V$ and $a_1,\ldots,a_d\in K$, then
\[
 \sum_{j=1}^d a_ju_j=0
 \quad\Longrightarrow\quad
 \sum_{j=1}^d[t^\delta]a_j\,u_j=0\quad\text{for every }\delta.
\]
Independence makes every scalar coefficient zero. Thus the tensor map is
injective and its image is $E_\Gamma(V)=\Span_KV$.

\begin{proposition}[Finite total coefficient rank]\label{prop:algebraic}
A vector $x\in V_\Gamma$ belongs to $E_\Gamma(V)$ if and only if all
its coefficients span a finite-dimensional $k$-subspace of $V$.
\end{proposition}
\begin{proof}
The coefficients of a finite sum $\sum a_ju_j$ lie in the span of its
finitely many $u_j$. Conversely, let $u_1,\ldots,u_d$ be a basis of the
coefficient span and write $x_\delta=\sum_j c_{j\delta}u_j$.
Each scalar series $a_j=\sum_\delta c_{j\delta}t^\delta$ has support
contained in $\supp x$, hence is a Hahn scalar. Then
$x=\sum_{j=1}^da_ju_j$.
\end{proof}

Finite total coefficient rank is much weaker than finite support. A
single constant vector multiplied by an arbitrarily long Hahn scalar
series still lies in $E_\Gamma(V)$.

\subsection{Exact characterization of the closure}
For a cut level $\gamma\in\Gamma$ define
\[
 V_{x,<\gamma}=\Span_k\{x_\delta:\delta<\gamma\}.
\]
The coefficients may have a transfinite well-ordered support, but this
is an ordinary set-generated vector subspace.

\begin{theorem}[Finite rank below every cut]\label{thm:completion}
For $\Gamma\ne0$,
\begin{equation}\label{eq:completion}
 C_\Gamma(V)=\{x\in V_\Gamma:
             \dim_k V_{x,<\gamma}<\infty\text{ for every }\gamma\in\Gamma\}.
\end{equation}
This is the valuation completion of $K\otimes_kV$ within $V_\Gamma$.
It is a $K$-subspace. Every vector of $V_\Gamma$ is a strong sum of
vectors of $E_\Gamma(V)$, whether or not it lies in $C_\Gamma(V)$.
\end{theorem}
\begin{proof}
Suppose $x$ lies in the closure. For any $\gamma$, choose
$y\in E_\Gamma(V)$ with $v(x-y)\ge\gamma$. Their coefficients below
$\gamma$ agree, and by \cref{prop:algebraic} those of $y$ have finite
span. Hence $V_{x,<\gamma}$ is finite-dimensional.

Conversely, assume the right-hand condition. For any $\gamma$, the
truncation
\[
 x_{<\gamma}=\sH_{\delta<\gamma}x_\delta t^\delta
\]
has finite total coefficient rank. It belongs to $E_\Gamma(V)$ and
satisfies $v(x-x_{<\gamma})\ge\gamma$. Raising $\gamma$ by a positive
amount proves approximation in every open valuation neighborhood.

A closure of a vector subspace in a topological vector space is a
subspace. Alternatively, addition combines two finite coefficient spans;
for multiplication by a fixed scalar $a$, coefficients below $\gamma$
use coefficients of $x$ below $\gamma-v(a)$, so their span is still
finite. The ambient space is complete by \cref{prop:spherical} below;
its closed subspace is therefore a completion of the embedded tensor
product. Finally $x=\sH_\delta t^\delta x_\delta$ is always a strong
sum of elements of $E_\Gamma(V)$.
\end{proof}

\begin{example}[A bounded increasing support]
\label{ex:bounded-support}
Let $e_2,e_3,\ldots$ be linearly independent and let $\Gamma$ be nonzero
and divisible. Choose $\alpha>0$ and put
\begin{equation}\label{eq:bounded-example}
 x=\sH_{n\ge2}t^{\alpha(1-1/n)}e_n.
\end{equation}
Its support is a strictly increasing sequence and is well ordered.
Every exponent is below $\alpha$, while the coefficients below that
cut have infinite rank. Consequently $x\notin C_\Gamma(V)$.
All finite partial sums belong to $E_\Gamma(V)$, but none approximates
$x$ with error of valuation at least $\alpha$. Nor can any other vector
of finite total coefficient rank do so.
\end{example}

The final sentence is essential. Showing only that the displayed
partial sums fail to converge would not exclude a different family of
approximants. The rank criterion excludes \emph{every} algebraic
scalar-extension approximant.

\subsection{Exactly which exponent groups make the algebraic part dense?}
\begin{lemma}[An order-theoretic alternative]\label{lem:group}
For a nonzero ordered abelian group $\Gamma$, precisely one of the
following holds: $\Gamma\simeq_{\mathrm{ord}}\Z$, or $\Gamma$ contains
a strictly increasing sequence bounded above in $\Gamma$.
Equivalently, every well-ordered subset has finite intersection with
every bounded initial interval exactly in the first case.
\end{lemma}
\begin{proof}
If $\Gamma$ has no least positive element, choose
$\delta_0>\delta_1>\cdots>0$. The sequence
$\delta_0-\delta_n$ for $n\ge1$ is strictly increasing and bounded
above by $\delta_0$.

Suppose instead that $\varepsilon$ is the least positive element.
If $\Gamma=\Z\varepsilon$, the asserted cyclic description holds.
Otherwise choose $\beta>0$ outside $\Z\varepsilon$. It must be larger
than every $n\varepsilon$. For if $n$ were the least positive integer
with $n\varepsilon>\beta$, then
$0<\beta-(n-1)\varepsilon<\varepsilon$, a contradiction.
Equality $\beta=n\varepsilon$ is excluded. Thus $n\varepsilon$ is a
bounded increasing sequence.

In an ordered copy of $\Z$, a nonempty well-ordered subset has a least
integer and therefore finitely many elements below any fixed integer.
Conversely, the increasing bounded sequence just constructed is itself
a well-ordered subset violating that property.
\end{proof}

\begin{theorem}[The cyclic-group boundary]\label{thm:cyclic}
If $\dim_kV<\infty$, then $E_\Gamma(V)=V_\Gamma$ for every $\Gamma$.
If $\dim_kV=\infty$ and $\Gamma\ne0$, then
\[
 C_\Gamma(V)=V_\Gamma\quad\Longleftrightarrow\quad\cyc.
\]
For infinite-dimensional $V$ and any nonzero $\Gamma$, the first
inclusion $E_\Gamma(V)\subset C_\Gamma(V)$ is strict whenever there
exists a sequence $n\alpha$ cofinal in $\Gamma$, with $\alpha>0$.
\end{theorem}
\begin{proof}
Finite dimension follows from \cref{prop:algebraic}. In the cyclic case,
every support has finitely many exponents below each cut; apply
\cref{thm:completion}. In the noncyclic case take a bounded increasing
sequence $\gamma_n<\beta$ from \cref{lem:group}, choose linearly
independent $e_n$, and form $\sum t^{\gamma_n}e_n$.
It violates the finite-rank condition below $\beta$.

For the last assertion, $x=\sum_{n\ge1}t^{n\alpha}e_n$ has infinite
total coefficient rank but only finitely many support terms below each
cut, by the stated cofinality. Hence it belongs to
$C_\Gamma(V)\setminus E_\Gamma(V)$.
\end{proof}

The last clause uses a sufficient condition that covers many familiar
workspaces. There is also an exact criterion for whether completion adds
\emph{any} new vectors.

\begin{theorem}[The cofinality boundary for nontrivial completion]
\label{thm:cofinality-completion}
Suppose $\dim_kV=\infty$ and $\Gamma\ne0$. Then
\[
 E_\Gamma(V)\subsetneq C_\Gamma(V)
 \quad\Longleftrightarrow\quad \operatorname{cf}(\Gamma)=\aleph_0,
\]
where $\operatorname{cf}(\Gamma)$ is the least cardinality of a cofinal
subset. Consequently the complete diagram is
\[
 \begin{array}{c|c}
 \text{Exponent group}&\text{Algebraic, completed, and full Hahn spaces}\\ \hline
 \Gamma\simeq_{\mathrm{ord}}\Z&E_\Gamma(V)\subsetneq C_\Gamma(V)=V_\Gamma\\
 \Gamma\text{ noncyclic},\ \operatorname{cf}(\Gamma)=\aleph_0
 &E_\Gamma(V)\subsetneq C_\Gamma(V)\subsetneq V_\Gamma\\
 \operatorname{cf}(\Gamma)>\aleph_0&E_\Gamma(V)=C_\Gamma(V)\subsetneq V_\Gamma.
 \end{array}
\]
\end{theorem}
\begin{proof}
A nonzero ordered group has no largest element, so its cofinality is
infinite. If it is countable, choose a strictly increasing cofinal
sequence $\gamma_n$ and independent coefficients $e_n$.
The vector $x=\sum_{n\ge1}t^{\gamma_n}e_n$ has finite coefficient rank
below each cut, since only finitely many $\gamma_n$ lie below that cut.
Its total coefficient rank is infinite. Thus it belongs to
$C_\Gamma(V)\setminus E_\Gamma(V)$.

Conversely, suppose the group has uncountable cofinality and
$x\in C_\Gamma(V)$. If its total coefficient rank were infinite,
choose countably many linearly independent coefficients of $x$.
Their exponents form a countable subset of $\Gamma$, which is not
cofinal. All are therefore below some common cut, contradicting the
finite-rank condition for that cut. So $x\in E_\Gamma(V)$.
Combine this with \cref{thm:cyclic} to obtain the displayed diagram.
\end{proof}

Thus an algebraic scalar extension can already be valuation-complete
while remaining strictly smaller than the full Hahn space. Completion
and closure under strong sums are genuinely different operations.
For $\Gamma=0$, the entire construction simply reduces to $V$ over $k$.

\section{Completeness and extension: the classical mechanism}
\label{sec:extension}

\subsection{Spherical completeness without a norm on the coefficients}
The extension argument needs full Hahn series, not merely a field of
left-finite or bounded-denominator series. The following proof makes
this distinction transparent.

\begin{proposition}[Spherical completeness of the full Hahn space]
\label{prop:spherical}
Every nested set-indexed family of nonempty closed valuation balls in
$V_\Gamma$ has nonempty intersection. In particular, $K$ is spherically
complete and $V_\Gamma$ is complete for its valuation uniformity.
No completeness or topology on $V$ is required.
\end{proposition}
\begin{proof}
Write the nested family as $B(a_j,\rho_j)$. Membership in that ball
means equality with $a_j$ at every exponent $\delta<\rho_j$.
For any exponent below at least one radius, define $x_\delta$ to be
that prescribed coefficient. Nested balls have compatible prescriptions:
their centers agree below the smaller radius. At exponents below no
radius, put $x_\delta=0$.

We must verify that $x$ is a Hahn vector, not just a function.
Let $A$ be a nonempty subset of its support and select $\delta_0\in A$.
A nonzero coefficient at $\delta_0$ was prescribed by some radius
$\rho_j>\delta_0$. All coefficients below $\rho_j$ agree with those
of $a_j$. The nonempty set $A\cap(-\infty,\rho_j)$ is therefore a
subset of a well-ordered support and has a least element. That element
is also the least of $A$, since every element not in the intersection
is larger. Thus $\supp x$ is well ordered. By construction $x$ is in
every ball.

For completeness one can argue directly as well. A Cauchy net has, for
each cut $\gamma$, eventually constant coefficients below $\gamma$.
These stabilized truncations are compatible. Assemble their coefficients
as above. The same well-ordering argument applies, and the resulting
Hahn vector is the limit, because every required truncation eventually
agrees. The zero coefficient space is immediate.
\end{proof}

\begin{remark}
In particular, spherical completeness here does not recover the ordinary
Hilbert topology on constants. Even an incomplete ordinary normed
coefficient space gives a complete Hahn space for this valuation. The
completeness is completeness of formal truncations.
\end{remark}

\subsection{A valuation Hahn--Banach lemma}
We record the exact extension statement used, with a proof that works
for arbitrary ordered groups. The rank-one version is classical
Ingleton theory \cite{Ingleton,Morillon}; the argument is the same
nested-ball argument, written additively.

\begin{theorem}[Contractive scalar extension]\label{thm:HB}
Let $X$ be a $K$-vector space with a valuation taking nonzero values in
$\Gamma$, satisfying the scalar and ultrametric laws. Let $M\le X$ and
$f:M\to K$ be $K$-linear with $v(fm)\ge v(m)$ for every $m\in M$.
There is an extension $L:X\to K$ with the same inequality.
\end{theorem}
\begin{proof}
First extend over one vector $y\notin M$. For each $m\in M$ consider
\[
 B_m=B\bigl(f(m),v(y-m)\bigr)\subseteq K.
\]
These balls form a chain. Indeed,
\[
 v(f(m)-f(n))\ge v(m-n)
 \ge\min\{v(y-m),v(y-n)\}.
\]
Thus the two centers lie within the ball of smaller valuation radius,
which is the larger ball, and the balls are nested. Their radii are
finite because $y\notin M$. By spherical completeness select
$c\in\bigcap_{m\in M}B_m$. Define
$g(m+ay)=f(m)+ac$.

The decomposition is unique. For $a\ne0$, put $n=-m/a$.
Since $c\in B_n$,
\[
 v(g(m+ay))=v(a)+v(c-f(n))
 \ge v(a)+v(y-n)=v(m+ay).
\]
The case $a=0$ is the original bound.

Finally order all contractive extensions of $f$ by extension. The union
along a chain is a well-defined contractive linear extension, so Zorn's
lemma supplies a maximal one. If its domain were not $X$, the one-step
argument would enlarge it. Thus its domain is all of $X$.
\end{proof}

The proof uses choice twice: to select an element of an intersection
when extending, and in the maximal-extension argument. No effective
procedure for an arbitrary input space is asserted. In particular, the
continuous functionals constructed from this theorem below are exact
existence results, not computable symbolic formulas on all Hahn vectors.

\section{The continuous dual and its invisible kernel}
\label{sec:continuous}

\subsection{The restriction space}
Let
\[
 \Dcont(V_\Gamma)=\Hom_K^{\mathrm{cont}}(V_\Gamma,K),
 \qquad
 \Dstr(V_\Gamma)=\Hom_K^{\mathrm{str}}(V_\Gamma,K).
\]
Let $V^\#=\Hom_k(V,k)$ be the ordinary \emph{algebraic} dual, not a
topological dual. By \cref{thm:strong},
\begin{equation}\label{eq:strong-dual}
 \Dstr(V_\Gamma)\simeq V^\#((t^\Gamma)).
\end{equation}

\begin{definition}\label{def:restriction-space}
The restriction space $P_\Gamma(V)$ consists of the $k$-linear maps
$A:V\to K$ such that there exists $s\in\Gamma$ with
$v(Au)\ge s$ for every $u\in V$. Its $K$-vector-space structure is
pointwise scalar multiplication.
\end{definition}

An equivalent coefficient description is a family
$(\varphi_\gamma)_{\gamma\in\Gamma}$ in $V^\#$ such that
$\{\gamma:\varphi_\gamma(u)\ne0\}$ is well ordered for each $u$, and
$\bigcup_u\{\gamma:\varphi_\gamma(u)\ne0\}$ is bounded below. The
value of $A(u)$ is the associated scalar Hahn series. In contrast,
$V^\#((t^\Gamma))$ requires this union itself to be well ordered.
The two requirements agree for a cyclic group, but not in general.

\begin{lemma}[Extension from constants to their algebraic span]
\label{lem:E-extension}
Every $A\in P_\Gamma(V)$ has a unique $K$-linear extension
$f_A:E_\Gamma(V)\to K$. If $v(Au)\ge s$ for all $u$, then
$v(f_Ae)\ge v(e)+s$ for all $e\in E_\Gamma(V)$.
\end{lemma}
\begin{proof}
Uniqueness and well-definedness follow from
$E_\Gamma(V)\simeq K\otimes_kV$. To check the bound, choose linearly
independent constants $u_1,\ldots,u_d$ with
$e=\sum_ja_ju_j$. Independence gives
$v(e)=\min_jv(a_j)$ after omitting zero scalars. Therefore
\[
 v(f_Ae)=v\left(\sum_ja_jA(u_j)\right)
 \ge\min_j\{v(a_j)+s\}=v(e)+s.
\]
The zero vector is automatic.
\end{proof}

\begin{theorem}[Exact restriction sequence]\label{thm:restriction}
Define
\[
 N_\Gamma(V)=\{L\in\Dcont(V_\Gamma):L(u)=0\text{ for every }u\in V\}.
\]
Restriction to constants yields an exact sequence of $K$-vector spaces
\begin{equation}\label{eq:exact-sequence}
 0\longrightarrow N_\Gamma(V)\longrightarrow\Dcont(V_\Gamma)
 \xrightarrow{\ \res\ }P_\Gamma(V)\longrightarrow0.
\end{equation}
Restriction is injective on $\Dstr(V_\Gamma)$; its image there is
exactly the embedded copy of $V^\#((t^\Gamma))$.
\end{theorem}
\begin{proof}
A continuous functional has a uniform valuation-shift bound by
\cref{lem:continuous-shift}. On nonzero constants the input valuation
is zero, so its restriction belongs to $P_\Gamma(V)$. The kernel is
exactly $N_\Gamma(V)$ by definition.

For surjectivity, let $A\in P_\Gamma(V)$ and take a lower bound $s$.
By \cref{lem:E-extension}, $t^{-s}f_A$ is contractive on
$E_\Gamma(V)$. Extend it to $V_\Gamma$ by \cref{thm:HB} and multiply
the extension by $t^s$. The result is continuous and restricts to $A$.
Finally, the assertion about strong maps is exactly
\cref{thm:strong}, with $W=k$.
\end{proof}

As vector spaces, the exact sequence admits a noncanonical splitting
by choice. There is no specified canonical splitting and no assertion
of a preferred extension operation. The sequence describes all
restrictions and all ambiguities, rather than concealing the ambiguity
in a purported coefficient formula for every continuous functional.

\subsection{Invisible functionals detect exactly the completion defect}
We call members of $N_\Gamma(V)$ \emph{invisible on constants}.
This is a definition, not a statement that they are zero: the constants
need not be topologically total in the Hahn space.

\begin{theorem}[Separation of the completion defect]
\label{thm:separation}
For every $V$ and every nonzero $\Gamma$,
\begin{equation}\label{eq:common-kernel}
 C_\Gamma(V)=\bigcap_{L\in N_\Gamma(V)}\ker L.
\end{equation}
More explicitly, if $x\notin C_\Gamma(V)$ and
$\dim_kV_{x,<\gamma}=\infty$, there is a $K$-linear functional $L$
satisfying
\begin{equation}\label{eq:ghost-values}
 L|_V=0,\qquad L(x)=t^\gamma,\qquad
 v(Ly)\ge v(y)\quad(y\in V_\Gamma).
\end{equation}
Such a functional is continuous but not strong.
\end{theorem}
\begin{proof}
Every member of $N_\Gamma(V)$ vanishes on the algebraic span of the
constants and, by continuity, on its closure. This proves one inclusion
in \eqref{eq:common-kernel}.

For the other, choose $x$ and $\gamma$ as in the statement. No
$e\in E_\Gamma(V)$ can agree with $x$ below $\gamma$, by the
coefficient-rank obstruction. Thus
\begin{equation}\label{eq:uniform-separation}
 v(x-e)<\gamma\qquad(e\in E_\Gamma(V)).
\end{equation}
Since $x\notin E_\Gamma(V)$, the formula
\[
 f(e+ax)=a t^\gamma
\]
is a well-defined linear map on $E_\Gamma(V)+Kx$.
If $a\ne0$, then
\[
 v(e+ax)=v(a)+v(x+a^{-1}e)<v(a)+\gamma=v(f(e+ax)).
\]
If $a=0$, the output is zero. Hence $f$ is contractive in valuation.
Apply \cref{thm:HB} to extend it to $L$ on all of $V_\Gamma$.
This proves \eqref{eq:ghost-values} and separates $x$.

Finally $x=\sH_\delta t^\delta x_\delta$, while
$L(t^\delta x_\delta)=0$ for every monomial term. Those images are a
summable zero family, but their sum is not $L(x)$. Therefore $L$ is
not strong.
\end{proof}

\begin{corollary}[Exact existence boundary]\label{cor:ghost-boundary}
For nonzero $\Gamma$, invisible nonzero continuous functionals exist
if and only if $V$ is infinite-dimensional and $\Gamma$ is not ordered
isomorphic to $\Z$. Moreover
$N_\Gamma(V)\cap\Dstr(V_\Gamma)=\{0\}$.
\end{corollary}
\begin{proof}
By \cref{thm:separation}, the kernel is nonzero exactly when the
algebraic scalar extension is not dense. Apply \cref{thm:cyclic}.
A strong functional vanishing on constants vanishes on every monomial
coefficient expansion, proving the intersection statement.
\end{proof}

\section{Two independent ways continuity can fail strongness}
\label{sec:twofailures}

\subsection{One failure changes sums; the other destroys summability}
The invisible functionals just constructed annihilate every term of a
particular strong sum but not its value. A different construction
produces an image family that is not even summable.

\begin{theorem}[Two failures, for every noncyclic group]
\label{thm:twofailures}
Assume $\dim_kV=\infty$ and $\Gamma\ne0$ is not ordered isomorphic
to $\Z$.
There exist continuous $K$-linear functionals of both following types.

\textup{(a)} A nonzero functional is zero on all constants and fails
commutation with a strong monomial sum whose termwise images are all
zero.

\textup{(b)} A functional has a pointwise finite-support restriction to
constants, bounded below in valuation, but sends some strongly summable
monomial family to the constant family $(1)_{n\ge1}$.
It cannot be written as a strong functional plus an invisible one.
\end{theorem}
\begin{proof}
Part (a) is \cref{thm:separation,cor:ghost-boundary}.
For (b), \cref{lem:group} provides an increasing sequence
$\gamma_n<\beta$. Put $\delta_n=\beta-\gamma_n$; these are strictly
decreasing and positive. Choose independent vectors $e_n$ and extend
them to a Hamel basis of $V$. Define $A:V\to K$ by
\[
 A(e_n)=t^{\delta_n},\qquad A(b)=0
 \quad\text{on the other basis vectors}.
\]
Every vector is a finite Hamel combination, so $A(u)$ has finite
support. All those supports are positive, so $A\in P_\Gamma(V)$
with lower bound zero. By \cref{thm:restriction}, extend $A$ to a
continuous $K$-linear functional $L$ on $V_\Gamma$.

The family $t^{-\delta_n}e_n$ is strongly summable, its exponents being
strictly increasing. But
\[
 L(t^{-\delta_n}e_n)=t^{-\delta_n}A(e_n)=1
 \qquad(n\ge1).
\]
The image family violates the finite-fiber condition at zero.
If $L=S+N$ with $S$ strong and $N$ invisible on constants, then
$S|_V=A$. This is impossible by \cref{thm:strong}, because the global
coefficient support of $A$ contains the strictly decreasing sequence
$\delta_n$.
\end{proof}

\begin{corollary}[A second quotient obstruction]\label{cor:second-quotient}
Restriction induces a canonical isomorphism
\begin{equation}\label{eq:second-quotient}
 \frac{\Dcont(V_\Gamma)}{\Dstr(V_\Gamma)+N_\Gamma(V)}
 \simeq
 \frac{P_\Gamma(V)}{V^\#((t^\Gamma))}.
\end{equation}
Under the hypotheses of \cref{thm:twofailures}, this quotient is
nonzero, as is $N_\Gamma(V)$ itself.
\end{corollary}
\begin{proof}
Apply the first and third isomorphism theorems to
\eqref{eq:exact-sequence}, using its identified image of the strong
subspace. Nonvanishing is the example in part (b).
\end{proof}

In part (a), preservation of summability is asserted only for the
displayed zero image family, not for every possible strong family.
The two examples give distinct obstructions; they do not assert logical
independence of the two universally quantified clauses in the definition.

Thus the full continuous dual is not obtained simply by adding
invisible functionals to the strong dual. There are also continuous
restrictions whose supports cannot be organized as one Hahn series of
coefficient functionals.

\subsection{Why ordinary Laurent exponents are exceptional}
\begin{theorem}[When continuous and strong duals agree]
\label{thm:continuous-strong}
For $\Gamma\ne0$,
\[
 \Dcont(V_\Gamma)=\Dstr(V_\Gamma)
 \quad\Longleftrightarrow\quad
 \bigl(\dim_kV<\infty\ \text{or}\ \cyc\bigr).
\]
\end{theorem}
\begin{proof}
If $V$ is finite-dimensional, its constants are a finite $K$-basis of
$V_\Gamma$. Any linear functional is represented by finitely many
scalar Hahn series, hence by a globally well-ordered coefficient
family, so is strong.

If $\Gamma\simeq_{\mathrm{ord}}\Z$, then $C_\Gamma(V)=V_\Gamma$,
so $N_\Gamma(V)=0$. For any $A\in P_\Gamma(V)$, its global support
is bounded below in an ordered copy of $\Z$, and is therefore well
ordered. Hence $P_\Gamma(V)=V^\#((t^\Gamma))$. The exact sequence
\eqref{eq:exact-sequence} now identifies every continuous functional
with the unique strong functional having its restriction.

In every remaining case, \cref{thm:twofailures} supplies continuous
nonstrong functionals.
\end{proof}

There is a useful topological explanation in the cyclic case. A strong
family has finitely many nonzero contributions below each cut, so its
finite subsums converge in valuation. A uniform valuation-shift bound
makes the image family have the same finiteness property after shifting
cuts. Continuity then forces commutation with the limit. For a
noncyclic group, a bounded increasing sequence is precisely the
obstruction to that argument.

\section{Hilbert coefficients: represented, strong, and bounded duals}
\label{sec:hilbert}

\subsection{The positive form and its field-valued norm}
Take $k=\C$, an ordinary complex Hilbert space $H$, and
\[
 K=\C((t^\Gamma)),\qquad F=\R((t^\Gamma)),\qquad
 \HG=H((t^\Gamma)).
\]
Order $F$ by its first nonzero coefficient and conjugate $K$
coefficientwise. Our inner products are linear in the first variable.
The coefficient-convolution form, already used in the repository
\cite{RepoSpectral}, is
\begin{equation}\label{eq:inner}
 \ip{x}{y}=\sH_{\alpha,\beta}
       \ip{x_\alpha}{y_\beta}_H t^{\alpha+\beta}.
\end{equation}
Each ordinary Hilbert inner product is evaluated as one coefficient
\emph{before} applying Hahn convolution. There is no strong summation
over infinitely many ordinary coordinates at the same exponent.

\begin{proposition}[Norm and valuation; standard coefficient extension]
\label{prop:norm}
The form \eqref{eq:inner} is positive definite. Its positive square-root
norm $\norm{x}=\sqrt{\ip{x}{x}}$ exists in $F$ for every ordered group
$\Gamma$, satisfies Cauchy--Schwarz and the triangle inequality, and
\[
 v(\norm{x})=v(x)\qquad(x\ne0).
\]
If $\delta=v(x)$, its leading coefficient is $\norm{x_\delta}_H$.
\end{proposition}
\begin{proof}
Convolution is defined by \cref{lem:support}. If $x\ne0$, the first
coefficient of $\ip{x}{x}$ is $\norm{x_\delta}_H^2>0$ at exponent
$2\delta$. Write it as
$c^2t^{2\delta}(1+u)$ with $c>0$ and $v(u)>0$, allowing $u=0$.
Then
\[
 \norm{x}=ct^\delta\sH_{n\ge0}\binom{1/2}{n}u^n
\]
is defined by the Hahn--Neumann lemma and squares to the form value.
Its first coefficient and exponent are as asserted. Notice that
$\delta$ is already in $\Gamma$: divisibility of $\Gamma$ is not
needed for these particular square roots. The case $H=\C$ defines
$|a|=\sqrt{a\bar a}$ for scalars as well.

Completing the square in
$\ip{x-ay}{x-ay}\ge0$, with
$a=\ip{x}{y}/\ip{y}{y}$ for $y\ne0$, proves
$|\ip{x}{y}|^2\le\ip{x}{x}\ip{y}{y}$.
Taking positive square roots gives Cauchy--Schwarz. Expansion of
$\norm{x+y}^2$, followed by that inequality, gives the triangle
inequality. These are finite algebraic manipulations over the ordered
field $F$.
\end{proof}

For $\Gamma\ne0$, the topology defined by all positive $F$-valued norm tolerances agrees
with the intrinsic valuation topology. For a given positive tolerance,
a strictly higher monomial valuation is smaller than that tolerance;
conversely, a suitably smaller monomial norm bound enforces any desired
valuation cutoff. This statement concerns the \emph{fixed} field $F$,
not the topology induced from the full surreal class.

\subsection{Boundedness is much weaker than ordinary coefficient boundedness}
\begin{proposition}[Order bounds for all continuous maps]
\label{prop:order-bound}
For $\Gamma\ne0$ and ordinary Hilbert spaces $H,G$, a $K$-linear map
$T:H_\Gamma\to G_\Gamma$ is valuation-continuous if and only if
there is $M\in F_{>0}$ such that
\begin{equation}\label{eq:order-bound}
 \norm{Tx}\le M\norm{x}\qquad(x\in H_\Gamma).
\end{equation}
In particular, every strong map is bounded in this sense, even if its
coefficient maps are ordinary discontinuous algebraic maps.
If $T$ is strong with least operator exponent $s$, then for each
$\eta>0$ the scalar $M=t^{s-\eta}$ is a valid bound.
\end{proposition}
\begin{proof}
Continuity gives $v(Tx)\ge v(x)+s$ for some $s$ by
\cref{lem:continuous-shift}. The valuation of
$t^{s-\eta}\norm{x}$ is strictly smaller than that of $\norm{Tx}$
when both are nonzero. In an ordered Hahn field this means the former
positive element is larger, irrespective of their ordinary leading
coefficients. The zero cases are immediate.
Conversely, a positive inequality \eqref{eq:order-bound} implies
$v(Tx)\ge v(M)+v(x)$: otherwise its left side would dominate its right
side. Apply \cref{lem:continuous-shift}. For strong maps use
\cref{cor:strong-calculus}.
\end{proof}

We do not define a real-Hahn operator norm by taking the least bound.
The set of bounds need not have a least element. The following theorem
specifies its exact form in a useful case.

\begin{theorem}[An ordinary unbounded functional has precisely infinite bounds]
\label{thm:infinite-bounds}
Let $\varphi:H\to\C$ be an ordinary unbounded algebraic linear
functional, and let $\widehat\varphi$ be its coefficientwise Hahn
extension. For $M\in F_{>0}$,
\[
 |\widehat\varphi(x)|\le M\norm{x}\quad\text{for all }x
 \quad\Longleftrightarrow\quad v(M)<0.
\]
Thus the admissible bounds are exactly the positive infinite Hahn
scalars, and this set has no least element.
\end{theorem}
\begin{proof}
Coefficientwise extension has $v(\widehat\varphi x)\ge v(x)$.
Every positive $M$ with $v(M)<0$ therefore dominates the output/input
norm ratio, proving sufficiency.
If $v(M)\ge0$, the scalar is finite, so $M<c$ for some positive
ordinary real $c$. Testing the proposed bound on constant vectors would
give $|\varphi(u)|\le c\norm{u}_H$ for all $u$, contradicting
unboundedness. Finally $M/2$ is a smaller positive infinite bound
whenever $M$ is one.
\end{proof}

Such $\varphi$ exist for every infinite-dimensional Hilbert space.
Choose an orthonormal sequence $e_n$, prescribe $\varphi(e_n)=n$ on
its algebraic span, and extend the sequence to a Hamel basis, prescribing
zero on the added basis vectors. The basis extension uses choice. The
unit vectors $e_n$ prove unboundedness. No continuity of this ordinary
functional is asserted.

\subsection{Exactly which strong functionals have representers?}
Write $H^\#=\Hom_\C(H,\C)$ and let $H'\subset H^\#$ be the ordinary
continuous dual. Define
\[
 \Drep=\{L_y:x\mapsto\ip{x}{y}\ ;\ y\in\HG\}.
\]
This is a $K$-subspace of the scalar-valued functionals: with our
convention $aL_y=L_{\bar a y}$. The map $y\mapsto L_y$ is conjugate
linear, not linear. This convention does not affect the following
linear description in terms of coefficient functionals.

\begin{theorem}[Coefficientwise Riesz criterion and exact first gap]
\label{thm:riesz}
There are canonical identifications of spaces of functionals
\begin{align}
 \Dstr&\simeq H^\#((t^\Gamma)),\label{eq:hilbert-strong}\\
 \Drep&\simeq H'((t^\Gamma)),\label{eq:hilbert-rep}\\
 \Dstr/\Drep&\simeq (H^\#/H')((t^\Gamma)).\label{eq:riesz-quotient}
\end{align}
A strong functional is inner-product-represented exactly when every
coefficient functional in its global Hahn expansion is ordinarily
continuous. Its representer is unique.
\end{theorem}
\begin{proof}
The first identification is \cref{thm:strong}. If $L=L_y$, then testing
on constant $u$ gives
$[t^\gamma]L(u)=\ip{u}{y_\gamma}_H$, an ordinarily continuous
functional. Conversely, let
$L=\sH_\gamma\varphi_\gamma t^\gamma$ be strong with each
$\varphi_\gamma\in H'$. The ordinary Hilbert-space Riesz
representation theorem gives a unique $y_\gamma\in H$ satisfying
$\varphi_\gamma(u)=\ip{u}{y_\gamma}_H$.
A coefficient functional is nonzero exactly when its representer is
nonzero, so these $y_\gamma$ have the same well-ordered support.
Their Hahn vector $y$ satisfies $L(x)=\ip{x}{y}$ by convolution.
Nondegeneracy proves uniqueness.

For the quotient, apply coefficientwise the quotient map
$q:H^\#\to H^\#/H'$. It defines a $K$-linear map on Hahn spaces
whose kernel is exactly $H'((t^\Gamma))$. It is onto: choose a
$\C$-linear section of $q$ and lift every coefficient through that
section, without enlarging the support. The quotient map and resulting
isomorphism are canonical; the auxiliary section need not be.
\end{proof}

This coefficientwise Riesz criterion is consistent with the repository's
classification of adjointable operators \cite{RepoSpectral}. Indeed, a
represented functional has an everywhere-defined adjoint from $K$ to
$\HG$, and conversely. We do not claim to have discovered a general
failure of Riesz representation in non-Archimedean analysis. What the
formula adds is an explicit description of the entire first gap for
this Hahn coefficient model.

\begin{corollary}[The three-dual trichotomy]\label{cor:three-duals}
Assume $\Gamma\ne0$.
If $H$ is finite-dimensional, then
$\Drep=\Dstr=\Dcont$.
If $H$ is infinite-dimensional and $\cyc$, then
\[
 \Drep\subsetneq\Dstr=\Dcont.
\]
If $H$ is infinite-dimensional and $\Gamma$ is noncyclic, then
\begin{equation}\label{eq:strict-three}
 \Drep\subsetneq\Dstr\subsetneq\Dcont.
\end{equation}
Every functional in these three spaces is bounded by some positive
$F$-scalar.
\end{corollary}
\begin{proof}
In finite dimension every ordinary functional is continuous and every
$K$-linear functional has a finite coordinate representation.
In infinite dimension an unbounded algebraic functional supplies the
first strict inclusion by \cref{thm:riesz}. The second equality or
strict inclusion follows from \cref{thm:continuous-strong}.
Boundedness is \cref{prop:order-bound}.
\end{proof}

\subsection{A cardinal lower bound on the first gap}
This optional quantitative statement shows that the first gap is not
accounted for by a few exceptional discontinuous coefficients.
\begin{proposition}\label{prop:dimension}
If $H$ is infinite-dimensional and separable and
$\mathfrak c=2^{\aleph_0}$, then
\[
 \dim_K(\Dstr/\Drep)\ge 2^{\mathfrak c}.
\]
No equality is asserted for arbitrary $\Gamma$.
\end{proposition}
\begin{proof}
An infinite-dimensional separable complex Hilbert space has cardinality
$\mathfrak c$ and Hamel dimension $\mathfrak c$.
For the lower bound on that dimension, fix an orthonormal sequence and
use the vectors $u_a=\sum_{n\ge0}a^ne_n$, $0<a<1$.
Every finite selection is linearly independent, by the Vandermonde
determinant obtained from its first coordinates. The cardinal upper
bound follows from separability, since vectors are determined by
countably many complex coordinates in an orthonormal basis.

The algebraic dual consists of all assignments of complex values to a
Hamel basis, so has cardinality
$\mathfrak c^{\mathfrak c}=2^{\mathfrak c}$.
A vector space over a field of size $\mathfrak c$ has cardinality equal
to the maximum of its infinite dimension and $\mathfrak c$.
Consequently $\dim_\C H^\#=2^{\mathfrak c}$, while
$\dim_\C H'=\mathfrak c$ by Riesz. Thus
$\dim_\C(H^\#/H')=2^{\mathfrak c}$.
Independent constant coefficient vectors remain independent over $K$:
extracting the least nonzero scalar exponent in a finite putative
relation gives a nontrivial complex relation. Apply
\eqref{eq:riesz-quotient}.
\end{proof}

\section{Closed hyperplanes with no orthogonal geometry}
\label{sec:hyperplane}

\subsection{A strong bounded projection with no orthogonal complement}
Let $\varphi:H\to\C$ be an unbounded algebraic functional as above,
let $M=\ker\varphi$, and choose $u\in H$ with $\varphi(u)=1$.
The ordinary algebraic hyperplane $M$ is dense in $H$.
Indeed, if its closure were proper, its codimension-one algebraic
maximality would force $M$ itself to be closed. A linear functional
with closed codimension-one kernel is continuous: orthogonal projection
onto its one-dimensional complement gives its bounded representation.
This would contradict the choice of $\varphi$.

Put $M_\Gamma=M((t^\Gamma))\subset H_\Gamma$.
The ordinary density of $M$ and the valuation closedness of $M_\Gamma$
are compatible because the two topologies treat coefficients
completely differently.

\begin{theorem}[A closed, strongly complemented hyperplane with zero
orthogonal complement]\label{thm:hyperplane}
For every nonzero $\Gamma$ and every infinite-dimensional ordinary
complex Hilbert space $H$, the subspace $M_\Gamma$ above has all the
following properties. It is a valuation-closed $K$-hyperplane; it is
closed under strong sums; and
\[
 P(x)=x-u\widehat\varphi(x)
\]
is a strong, order-bounded projection onto it. Nevertheless
\[
 M_\Gamma^\perp=\{0\},
\]
so there is no orthogonal projection onto $M_\Gamma$ and
$H_\Gamma\ne M_\Gamma\oplus M_\Gamma^\perp$.
\end{theorem}
\begin{proof}
The coefficientwise map $\widehat\varphi$ is onto $K$, because
$\widehat\varphi(au)=a$. Its kernel is exactly $M_\Gamma$.
It is continuous, so that kernel is closed and has codimension one.
Directly, if $x\notin M_\Gamma$, choose a coefficient
$x_\delta\notin M$; every vector agreeing with $x$ through that
coefficient is outside $M_\Gamma$.
Strong sums stay in $M_\Gamma$ since every coefficient sum is finite
and lies in $M$.

The ordinary algebraic map $p=I-u\varphi$ has $p^2=p$ and range $M$.
Its coefficientwise Hahn extension is precisely $P$.
It is strong by \cref{thm:strong}, and bounded by
\cref{prop:order-bound}.

If $y\perp M_\Gamma$, testing against every constant $m\in M$
gives $\ip{m}{y_\delta}_H=0$ for every $\delta$.
Ordinary density of $M$ forces $y_\delta=0$ for all $\delta$, hence
$y=0$. A proper subspace cannot then have an orthogonal decomposition
with its orthogonal complement. An orthogonal projection would give
such a decomposition, so none exists.
\end{proof}

This is stronger than saying that a particular constructive formula for
an orthogonal projection fails. The subspace has codimension one and a
strong bounded algebraic projection, but no possible orthogonal
complement. The general non-Archimedean phenomenon has predecessors
\cite{AN}; the role of the coefficient topology is fully explicit here.

\subsection{The distance does not merely fail to attain a minimum}
\begin{theorem}[Exact lower-bound cut for the distance]
\label{thm:distance}
For a constant $u\notin M$, consider
\[
 \mathcal D_u=\{\norm{u-m}:m\in M_\Gamma\}\subset F_{>0}.
\]
Its nonnegative lower bounds are exactly zero and all positive
infinitesimals of $F$. Consequently $\mathcal D_u$ has no infimum in
$F$.
\end{theorem}
\begin{proof}
For $m\in M_\Gamma$, its coefficient $m_0$ lies in $M$, so
$u-m_0\ne0$. Thus $u-m$ cannot have positive valuation.
If $m$ has a negative-valuation leading term, then $\norm{u-m}$ is
positive infinite. Otherwise its norm has valuation zero with positive
ordinary leading coefficient. In either case every positive
infinitesimal is smaller than $\norm{u-m}$.

Conversely, density of $M$ in the ordinary Hilbert topology gives,
for every ordinary $r>0$, a \emph{constant} $m\in M$ with
$\norm{u-m}_H<r$. A nonnegative lower bound in $F$ must therefore be
smaller than every positive real. It is either zero or positive
infinitesimal. This identifies the lower-bound set exactly.
It has no greatest element: any positive infinitesimal can be doubled,
and there exists at least one positive infinitesimal because
$\Gamma\ne0$. Hence no infimum exists.
\end{proof}

The ordinary distance from $u$ to $M$ is zero. The Hahn field does not
replace that distance by a particular positive infinitesimal. Rather,
all positive infinitesimals lie below every available distance, and
there is no greatest one. Confusing an infinitesimal scale with a
complete ordered continuum would miss this distinction.

\section{Enlarging the exponent group can undo completion}
\label{sec:enlargement}

\subsection{The exact cofinality dichotomy}
Let $\Gamma\subseteq\Delta$ be an ordered subgroup inclusion, and
embed $K_\Gamma$ in $K_\Delta$ and $V_\Gamma$ in $V_\Delta$ by
preserving coefficients and exponents. These embeddings preserve
strong sums. Each group supplies its own intrinsic valuation topology.
Recall that $\Gamma$ is cofinal in $\Delta$ if for every
$\delta\in\Delta$ some $\gamma\in\Gamma$ satisfies
$\delta\le\gamma$.

\begin{theorem}[Cofinal preservation and noncofinal collapse]
\label{thm:basechange}
For nonzero $\Gamma\subseteq\Delta$ and any $V$,
\begin{equation}\label{eq:algebraic-intersection}
 E_\Delta(V)\cap V_\Gamma=E_\Gamma(V),
\end{equation}
and
\begin{equation}\label{eq:completion-basechange}
 C_\Delta(V)\cap V_\Gamma=
 \begin{cases}
 C_\Gamma(V),&\Gamma\text{ is cofinal in }\Delta,\\
 E_\Gamma(V),&\Gamma\text{ is not cofinal in }\Delta.
 \end{cases}
\end{equation}
These are intersections of explicitly embedded vector sets; no claim
that a completed tensor-product functor commutes with extension is
implicit.
\end{theorem}
\begin{proof}
The algebraic criterion is finite total coefficient rank, which is
unchanged by the embedding. This proves
\eqref{eq:algebraic-intersection}.

If $\Gamma$ is cofinal and $x\in C_\Gamma(V)$, then for every
$\delta\in\Delta$ choose $\gamma\in\Gamma$ with
$\gamma>\delta$; strictness is obtained by adding a positive group
element. The coefficients of $x$ below $\delta$ lie among those below
$\gamma$ and have finite rank. Thus $x\in C_\Delta(V)$.
The converse follows by testing only cuts in $\Gamma$.

If $\Gamma$ is not cofinal, some $\delta\in\Delta$ is larger than
every element of $\Gamma$. For $x\in V_\Gamma$, \emph{all} its
coefficients occur below this one cut. Membership in $C_\Delta(V)$
therefore forces finite total coefficient rank, or
$x\in E_\Gamma(V)$. The reverse inclusion is automatic because an
algebraic vector belongs to the completion in every larger workspace.
\end{proof}

\subsection{A completely explicit two-scale example}
Take $\Gamma=\Q$ and
$\Delta=\Q\oplus_{\mathrm{lex}}\Q$, embedding $q\mapsto(0,q)$.
Both groups are divisible. Let $e_n$ be linearly independent and form
\begin{equation}\label{eq:two-scale}
 x=\sH_{n\ge1}t^{(0,n)}e_n,
\end{equation}
viewed first as $\sum_{n\ge1}t^ne_n$ in $V_\Q$.
There are finitely many exponents below each rational cut, so
$x\in C_\Q(V)$, and its finite partial sums converge in that intrinsic
topology. Its total coefficient rank is infinite, so
$x\notin E_\Q(V)$.

In $\Delta$, the single exponent $\beta=(1,0)$ exceeds every
$(0,q)$. Below this new cut, the coefficients of $x$ have infinite
rank. Hence $x\notin C_\Delta(V)$. Every finite coefficient-span
approximant $y$, even using scalars from $K_\Delta$, has
\[
 v(x-y)<(1,0).
\]
Admitting the scale $t^{(1,0)}$ defeats \emph{all} such approximants,
not just the original partial sums.

In addition, \cref{thm:separation} supplies a continuous
$K_\Delta$-linear functional $L$ on $V_\Delta$ with
\[
 L|_V=0,\qquad L(x)=t^{(1,0)}.
\]
No nonzero functional with these values could be continuous in the
original topology on $V_\Q$, since there $x$ was a limit of vectors
in the algebraic constant span. There is no contradiction: the
new induced topology on the old vectors is finer. In fact it is
discrete, since the ball of valuation greater than $(1,0)$ intersects
$V_\Q$ only in zero.

\subsection{What remains coherent under extension}
The support expansion of a strong operator
$A\in\Hom_k(V,W)((t^\Gamma))$ is still a well-ordered expansion
after embedding $\Gamma$ into $\Delta$. It therefore defines a
strong $K_\Delta$-linear operator on $V_\Delta$, and its action on
$V_\Gamma$ agrees with the old action. Composition remains coherent.
The same is true for coefficientwise represented functionals, ordinary
unbounded algebraic functional extensions, and the hyperplane
construction.

By contrast, an invisible continuous functional comes from extension
choices on vectors not determined by constants. Its definition does
not supply a canonical coefficientwise enlargement, and the completion
on which it must vanish changes exactly as
\eqref{eq:completion-basechange} specifies. We claim no canonical
transport for such functionals.

\section{Interpretation inside the surreal and surcomplex numbers}
\label{sec:surreal}

\subsection{Fixed Hahn workspaces give genuine surreal scales}
The link to surreals is the normal-form embedding. For a set-sized
ordered additive subgroup $\Gamma\subset\No$, send
\begin{equation}\label{eq:surreal-embedding}
 t^\gamma\longmapsto\omega^{-\gamma}.
\end{equation}
A well-ordered support in $\Gamma$ becomes a reverse-well-ordered
support of Conway exponents. Normal forms therefore embed
$\R((t^\Gamma))$ as an ordered subfield of $\No$ and
$\C((t^\Gamma))$ into $\No[i]$. This standard correspondence is
reviewed in the modern surreal automorphism literature
\cite{KKS} and implemented in the repository's normal-form layer
\cite{RepoRoot}. For this application one may take divisible groups,
so that the scalar workspaces are also real closed and algebraically
closed, respectively; the main duality proofs did not require that
extra hypothesis.

The Hilbert coefficients themselves remain vectors in a fixed set-sized
ordinary space $H$. We are constructing a module over an actual
surcomplex subfield, not claiming that a Hilbert vector is a single
surcomplex scalar. Inner products, norms, operator values, and their
inequalities are honest scalars in those subfields and hence in the
surreal or surcomplex numbers.

\subsection{The two-scale example in Conway notation}
For the preceding example take
\[
 \Gamma=\Q\subset\No,\qquad
 \Delta=\Q\omega+\Q\subset\No,
\]
with the natural lexicographic ordering under
$(a,b)\mapsto a\omega+b$. Then \eqref{eq:two-scale} becomes
\begin{equation}\label{eq:omega-example}
 x=\sH_{n\ge1}\omega^{-n}e_n.
\end{equation}
It is a limit of finite sums for the intrinsic rational-exponent
valuation topology. After admitting the exponent $\omega$, no finite
coefficient-span vector approximates it to valuation $\omega$.
The new scalar tolerance is $\omega^{-\omega}$, and an invisible
functional in the enlarged workspace can satisfy
$L(x)=\omega^{-\omega}$ while vanishing on every constant vector.

The support of \eqref{eq:omega-example} and its strong sum have not
changed. What changed is the family of precision demands. This gives a
concrete obstruction to constructing infinite-dimensional surcomplex
analysis by completing finite-coefficient data at one stage and
assuming the completion is unaffected by larger surreal scales.

\subsection{Three scope restrictions}
All theorems above concern set-sized supports, exponent groups, vector
spaces, and index sets. They are ordinary ZFC statements. They do not
apply Zorn's lemma to a proper class of vectors, and they do not build a
continuous dual of a class-sized space over the whole of $\No[i]$.

The intrinsic topology on a fixed Hahn workspace must not be identified
with the subspace topology inherited from the full surreal fine
topology. A set-sized collection of surreal exponents admits still
larger surreal exponents, and those extra tolerances can isolate old
points. The repository already emphasizes the distinction between
strong summability, workspace convergence, and full-surreal fine
convergence \cite{RepoInventory}. Our completion-enlargement theorem
is an algebraically precise instance of that distinction.

Finally, coarse order-boundedness is not an ordinary operator norm
estimate with a real constant. A positive infinite bound hides all
ordinary finite coefficient growth. Extra hypotheses are required to
recover the ordinary analytic meaning of boundedness, adjoints, or
Riesz representation. The coefficientwise criterion
\cref{thm:riesz} states exactly one such hypothesis.

\section{Worked checks and a formalization route}\label{sec:checks}

\subsection{What the accompanying program does}
The accompanying \texttt{verify\_examples.py} uses exact rational
arithmetic and finite-dimensional rational linear algebra. It checks
finite coefficient convolution, composition and scalar compatibility,
leading exponents of squared norms, coefficient-rank tests, and finite
instances of the projection and two-scale formulas. Its output is
recorded in \texttt{verification\_report.json}.

The program does not claim to certify the infinite theorems.
Well-ordering arguments, spherical completeness, Hamel-basis extension,
and the existence of invisible functionals are proved in the text.
In particular, no finite truncation can witness an invisible functional
that is zero on constants but nonzero on a Hahn sum: a finite
truncation is a finite $K$-linear combination of constants and must be
annihilated. This limitation is itself a consequence of linearity,
not an implementation defect.

\subsection{A small finite calculation that does carry information}
Let $V=\Q^3$, $W=\Q^2$, and use two-dimensional rational lexicographic
exponents. A finite operator series
$A=A_0+t^{(0,1)}A_1+t^{(1,0)}A_2$ and finite vector series
$x=x_0+t^{(0,2)}x_1$ give
\begin{align*}
 Ax={}&A_0x_0+t^{(0,1)}A_1x_0+t^{(0,2)}A_0x_1
       +t^{(0,3)}A_1x_1\\
     &+t^{(1,0)}A_2x_0+t^{(1,2)}A_2x_1.
\end{align*}
The checks compare this coefficient-level action with the action of
composed operator series. Cancellations are retained rather than
assuming every possible support sum is occupied. Rank checks similarly
use actual linear dependence, not just the number of nonzero
coefficients. For example, arbitrarily many coefficients of the form
$c_nu$ have rank one and do not obstruct algebraic scalar extension.

\subsection{Dependencies for a Lean development}
A feasible formalization separates three layers. First comes the
purely algebraic Hahn-vector layer: well-ordered supports, finite fibers,
coefficient maps, convolution, and the descending-support test in
\cref{thm:strong}. This layer needs no Hilbert space. Second comes the
valuation layer: balls as matching truncations, coefficient-rank closure,
spherical completeness, and the one-step extension theorem followed by
Zorn. Third comes the Hilbert layer: the coefficient form, its binomial
norm, ordinary Riesz representation, and the exact dual comparison.

The repo already has substantial scalar Hahn infrastructure and an
independent normal-form bridge \cite{RepoRoot}. Those are reusable
prerequisites, not a claim that the new module theorems are already
formalized. A useful first milestone would be the finite coefficient-rank
criterion and the cofinal/noncofinal intersection theorem: they isolate
the key obstruction without first developing an entire operator theory.
The nonconstructive steps should be explicit uses of classical choice,
not hidden as executable definitions of arbitrary functionals.

\section{Limits, further questions, and the audit}\label{sec:audit}

\subsection{What the article settles and what remains separate}
Within the specified Hahn-vector model, the completion question,
strong-Hom question, image and kernel of continuous restriction,
existence of invisible functionals, and scalar-enlargement question are
settled by the displayed theorems. The Hilbert specialization determines
exactly which strong functionals are represented and exactly when the
three duals coincide.

Several separate problems remain worthwhile. One is to classify
additional regularity conditions on restrictions to constants that
recover a useful intermediate dual while remaining stable under
noncofinal scalar enlargement. A second is to describe the size and
possible canonical structures of $N_\Gamma(V)$ beyond its exact
annihilator characterization. A third is to develop operator ideals in
this coefficient model that retain ordinary compactness or nuclearity
without mistaking valuation continuity for an ordinary norm estimate.
These are proposed continuations, not assertions of named published
open problems or prerequisites for the results proved here.

\subsection{Repository comparison}
The repository was inspected at the fixed revision
\begin{center}\small\texttt{\Commit}\end{center}
The principal comparison sources were the root README, the complete
report inventory in \texttt{docs/README.md}, the introduction and
relevant Hahn--Hilbert construction/operator sections of
\texttt{docs/surcomplex/infinite-dimensional-hahn-spectral-theory/article.tex},
and the introductions to the differential-equations and
analytic-geometry reports. The latter two were examined to avoid
reusing directions already substantially developed in the project.
The report inventory provides the comparison map \cite{RepoInventory}.

The closest pre-existing result is the automatic adjoint theorem in the
infinite-dimensional spectral report. It identifies adjointable
endomorphisms with $\mathcal B(H)((t^\Gamma))$ and uses ordinary
closed-graph and Baire arguments. The present strong-Hom theorem
identifies \emph{strong} endomorphisms with the larger algebra
$\End_\C(H)((t^\Gamma))$ by a different, purely support-theoretic
argument. The principal additional results concern the completed
constant span and the continuous dual beyond that algebra. We found
no statement of the exact coefficient-rank completion theorem, the
restriction exact sequence with its two failure mechanisms, or the
cofinal/noncofinal completion-intersection formula in the inspected
material. This is a targeted content comparison, not an exhaustive
audit of every line or historical revision of the repository.

\begin{table}[ht]
\centering\small
\begin{tabularx}{\textwidth}{@{}p{.24\textwidth}YY@{}}
\toprule
Topic & Established input or nearby result & Role in this article\\
\midrule
Strong summability & Hahn theory; BKKPS summability spaces & Definitions and support machinery, not new\\
Hahn--Banach extension & Ingleton; Morillon's choice discussion & Full needed proof given; classical mechanism\\
Hahn--Hilbert adjoints & Repository spectral article & Distinguished from strong and continuous maps\\
Completion & Algebraic scalar extension inside a Hahn module & Exact finite-rank-below-cuts criterion and cyclic boundary\\
Continuous restriction & Uniform valuation estimates and extension & Exact image, invisible kernel, two independent obstructions\\
Exponent enlargement & Coefficient-preserving Hahn embeddings & Cofinal preservation versus noncofinal collapse\\
Riesz and orthogonality & Ordinary Riesz; known non-Archimedean failures & Exact quotient and explicit hyperplane/distance-cut package\\
\bottomrule
\end{tabularx}
\caption{Attribution and the proposed contribution. Closely related
classical tools are not recast as new theorems.}
\end{table}

\subsection{Literature and verification boundary}
The literature check consulted the primary summability article
\cite{BKKPS}, the surreal automorphism article \cite{KKS}, Morillon's
paper \cite{Morillon}, and publication records for the neighboring
non-Archimedean Hilbert literature \cite{AN}. Bibliographic details of
Ingleton's original paper are corroborated in Morillon's references;
we did not obtain its full text. Some publisher pages exposed only
metadata. Search results that were irrelevant or inaccessible were not
treated as evidence of absence. The references below identify the
actual theoretical predecessors rather than trying to certify priority
by a keyword count.

The main arguments have been checked internally for their hypotheses,
including arbitrary rank, noncofinal inclusions, zero maps, finite
coefficient dimension, and the cyclic exceptional case. The exact
finite examples have been executed. Neither activity is independent
refereeing, and no Lean verification of these new statements is
claimed. The package is a research manuscript with complete proofs,
not a certified resolution of a longstanding named conjecture.

\section{Conclusion}
There are two independent reasons why extending a Hilbert space to Hahn
scalars does not transport the classical duality picture unchanged.
First, valuation boundedness ignores ordinary coefficient growth: an
ordinary discontinuous algebraic functional becomes strongly summable
and bounded by every positive infinite scalar. Second, valuation
completion does not exhaust strong Hahn summation: a vector with
infinitely many independent coefficients below one cut cannot be
reconstructed from finite coefficient spans, and continuous functionals
can detect precisely this defect while vanishing on every constant.

The organizing formulas are
\[
 \Hom_K^{\mathrm{str}}(V_\Gamma,W_\Gamma)
 =\Hom_k(V,W)((t^\Gamma))
\]
under convolution and
\[
 C_\Gamma(V)=
 \{x:\text{the coefficient rank below every cut is finite}\}.
\]
Their interaction yields the exact continuous restriction sequence,
the three-dual trichotomy, and a sharp failure of completion under
noncofinal enlargement. The concrete tolerance $\omega^{-\omega}$
in \cref{eq:omega-example} illustrates why these distinctions matter
specifically at surreal scales: adding a new scale can invalidate an
old completion without changing a single coefficient of its strong sum.

\begin{thebibliography}{99}
\small\raggedright\setlength{\itemsep}{.25em}
\bibitem{BKKPS}
V. Bagayoko, L. S. Krapp, S. Kuhlmann, D. Panazzolo, and M. Serra,
\emph{Automorphisms and derivations on algebras endowed with formal
infinite sums}, arXiv:2403.05827v2, September 22, 2025.
\url{https://arxiv.org/abs/2403.05827v2}.
See especially Sections 1.1--1.3.

\bibitem{Ingleton}
A. W. Ingleton, \emph{The Hahn--Banach theorem for non-Archimedean
valued fields}, Proceedings of the Cambridge Philosophical Society
\textbf{48} (1952), 41--45.
Bibliographic record and one-step statement also appear in
\cite{Morillon}, Section 3.1.2 and reference [9].

\bibitem{Morillon}
M. Morillon, \emph{Multiple Choices imply the Ingleton and Krein--Milman
axioms}, arXiv:1901.04021, 2019.
\url{https://arxiv.org/abs/1901.04021}.

\bibitem{AN}
J. Aguayo and M. Nova, \emph{Non-Archimedean Hilbert like spaces},
Bulletin of the Belgian Mathematical Society--Simon Stevin
\textbf{14} (2007), no. 5, 787--797.
\url{https://doi.org/10.36045/bbms/1197908895}.
The publication record and abstract were accessible; the complete
publisher PDF was not available to this inspection.

\bibitem{Neumann}
B. H. Neumann, \emph{On ordered division rings}, Transactions of the
American Mathematical Society \textbf{66} (1949), 202--252.
The classical ordered-series support lemma; also discussed in the
repository's Hahn support preliminaries \cite{RepoSpectral}.

\bibitem{KKS}
E. Kaplan, L. S. Krapp, and M. Serra,
\emph{Decomposing the automorphism group of the surreal numbers},
arXiv:2509.22374v3.
\url{https://arxiv.org/abs/2509.22374v3}.
Used for normal-form conventions and comparison with established
strong-linearity terminology, not as an input to the duality proofs.

\bibitem{RepoRoot}
V. Reshetnikov, \emph{Surreal}, repository, root README and normal-form
infrastructure, revision \texttt{\Commit}, inspected September 22, 2026.
\url{https://github.com/VladimirReshetnikov/Surreal/tree/048b72cf7cbfc8ab246e4f73788c10460cb3f6e0}.

\bibitem{RepoInventory}
\emph{The surreal and surcomplex reports}, repository report inventory,
\texttt{docs/README.md}, same revision as \cite{RepoRoot}.
\url{https://github.com/VladimirReshetnikov/Surreal/blob/048b72cf7cbfc8ab246e4f73788c10460cb3f6e0/docs/README.md}.

\bibitem{RepoSpectral}
\emph{Infinite-Dimensional Hahn Spectral Theory}, research manuscript
in the Surreal repository,
\texttt{docs/surcomplex/infinite-dimensional-hahn-spectral-theory/article.tex},
same revision as \cite{RepoRoot}.
\url{https://github.com/VladimirReshetnikov/Surreal/blob/048b72cf7cbfc8ab246e4f73788c10460cb3f6e0/docs/surcomplex/infinite-dimensional-hahn-spectral-theory/article.tex}.
See Part I: coefficient Hahn--Hilbert spaces, automatic adjoint
structure, and spectral range defects.
\end{thebibliography}
\end{document}
